\documentclass[10pt,a4paper]{amsart}
\usepackage[margin=19mm]{geometry}
\usepackage{amsmath,amssymb,mathtools}
\usepackage[scr=rsfs,cal=euler]{mathalfa}
\usepackage{xfrac}
\usepackage[unicode,hidelinks,bookmarks=false]{hyperref}
\newcommand{\ie}{i.e.\ }
\newcommand{\cf}{cf.\ }
\newcommand{\resp}{respectively\ }
\newcommand{\Subsection}{\subsection}
\providecommand{\nobreakdash}{\nobreak\hbox{-}}
\setlength{\emergencystretch}{2em}
\multlinegap0pt
\usepackage{bm}
\usepackage{bbm}
\usepackage{dsfont}
\usepackage{xspace}
\usepackage{cancel}
\usepackage{mathtools}
\usepackage{amsthm}
\makeatletter
\@ifundefined{theorem}{\newtheorem{theorem}{Theorem}}{}
\@ifundefined{prop}{\newtheorem{prop}[theorem]{Proposition}}{}
\makeatother
\newcommand{\cc}{\mathbb{C}}
\newcommand{\dif}{\mathrm{d}}
\newcommand{\zz}{\mathbb{Z}}
\title[The Schwinger-Dyson equations for random fuzzy geometries co]{The Schwinger-Dyson equations for random fuzzy geometries coupled to matter}
\author{Jeremy Gamble, Masoud Khalkhali, Nathan Pagliaroli}
\date{Source: arXiv:2606.01343v1 (CC BY 4.0)}
\begin{document}
\maketitle

\noindent\emph{Source and licence.} The Schwinger-Dyson equations for random fuzzy geometries coupled to matter. Version arXiv:2606.01343v1 (CC BY 4.0), distributed under the Creative Commons Attribution 4.0 International licence, as recorded in the arXiv licence metadata for this exact version and shown on the abstract page https://arxiv.org/abs/2606.01343v1. Full rights notice: licenses/arxiv-2606.01343-CC-BY-4.0.txt.\par\smallskip

\noindent\textit{Editorial scope.} Counted unit: the Prop whose source label is \texttt{None} in the article (source0.tex lines 329--331, source label \texttt{None}), delivered together with the article's own complete original proof of it (source0.tex lines 333--381, one \texttt{proof} environment of 49 lines). No mathematics is written, repaired or re-derived: statement and proof are the article's own text line for line. Required context retained uncounted: none. Every definition, display and label of those lines is retained unchanged. Omitted from this package and not redistributed: 1--328, 332--332, 382--1122. Nothing inside a retained excerpt is omitted or altered: every retained body is delivered exactly as the article's lines are written, so no omission receipt of any kind accompanies this package. Citations: the retained excerpts carry no citation command at all, so no bibliography entry of the article is transcribed and no reference is redistributed. Reference adaptation: none was needed and none was made; every reference inside the delivered unit resolves inside it, so no unresolved reference or citation is produced. Printed slips kept literal: no printed word, symbol, hypothesis or number of the counted statement or of its proof was altered. Layout and class: the article's own class is \texttt{article} with packages including tikz, authblk, graphicx, amsfonts, float, mathtools, subfigure, amssymb, latexsym, amsmath, amscd, slashed, mathrsfs, bm; the wrapper is the lane's pinned \texttt{amsart} editorial wrapper at 10pt on A4, which restates the theorem-like environments the retained text needs and adds the article's own macro definitions that the retained text needs (\texttt{\textbackslash cc}, \texttt{\textbackslash dif}, \texttt{\textbackslash zz}), with extra wrapper packages (bm, bbm, dsfont, xspace, cancel, mathtools). The delivered numbering is the unit's own. Assets: no retained span contains a figure environment, a graphics-inclusion command, a PDF-page inclusion, a table or a TikZ picture; no image file is copied into this package, the package declares no asset dependency and no image file is redistributed with it. Licence: version arXiv:2606.01343v1 of this article is distributed under the Creative Commons Attribution 4.0 International licence, as recorded in the arXiv licence metadata for this exact version and shown on the abstract page https://arxiv.org/abs/2606.01343v1. A full rights notice is delivered at licenses/arxiv-2606.01343-CC-BY-4.0.txt. Characters: every retained excerpt is pure ASCII, so the delivered engine, XeLaTeX with its default Latin Modern text font, reports no missing character.
\begin{prop}
	The function $H$ is entire. 
\end{prop}

\begin{proof}
	Before we prove that $F$ converges to an entire function, we first show that this function cannot have cuts along any $I_k$. Indeed, the function $f_k(x)$ has cuts along $I_k$ and $I_{k+1}$, however, observe
	\begin{align*}
		f_{k-1}(x)+f_k(x)+f_{k+1}(x) &=  g_k(x) + W(x_{k-1})^2 + W(x_{k+1})^2 -\frac{\beta_2}{2}W(x_{k+1})W(x_{k+2}) \\
		&\qquad - \frac{4}{\beta}[ V'(x_{k-1})W(x_{k-1})+V'(x_{k+1})W(x_{k+1})\\
		&\qquad - R(x_k)-R(x_{k-1})-R(x_{k+1}) ].
	\end{align*}
	Thus, $f_{k-1}+f_k + f_{k+1}$ is a function without cuts along $I_k$, since $g_k$ has no cuts along $I_k$, and $R$ is polynomial and therefore has no cuts at all. Considering that $f_k$ may only have cuts along $I_{k-1},I_{k},I_{k+1}$, and $F$ may be written as 
	\begin{align*}
		H(x) = f_{k}(x)+f_k(x)+f_{k+1}(x) + \sum_{j\neq k-1,k,k+1} f_j(x),
	\end{align*}
	i.e. a sum of two functions, neither of which has a cut along $I_k$. Since the choice of $k$ was arbitrary, we conclude that $S$ does not have cuts along any $I_k$. \\
	
	
	The previous argument immediately shows that $\sum_{|k|\leq N} f_k(x)$ has no cuts until $I_{\pm N}$, which will be useful in proving the convergence of $F$, which we turn to now. We will show that $H$ converges uniformly on all open balls $B(0;R) \coloneqq \{ x\in \cc : |x|<R \}$. \\
	
	We begin with some estimates that will prove useful. Recall that $L$ is the length of $\Sigma$. This gives a bound on the moments,
	\begin{align*}
		|\mu_n| \leq  \int_\Sigma \bigg| x^n \rho(x)\bigg| \dif x \leq L^n \int_\Sigma \rho(x)\dif x = L^n.
	\end{align*}
	
	Let us now assume $|x|<R$. Taking a term $W(x_k)$, we can expand this when 
	\[ L<|k|m-R<|k|m-|x| \leq \big||x|-|kmi|\big|\leq |x_k|  \]
	i.e. when $|k|>(R+L)/m$, in which case we have
	\begin{align*}
		|W(x_k)|&=\bigg|\sum_{n=0}^\infty \frac{\mu_n}{x_k^{n+1}}\bigg| \leq \frac{1}{|x_k|}\sum_{n=0}^\infty \bigg(\frac{L}{|x_k|}\bigg)^{n}  = \frac{1}{|x_k|-L} \leq \frac{1}{|k|m-(R+L)}.
	\end{align*}
	We now examine the term $V'(x_k)W(x_k)-R(x_k)$ carefully. We will write down the potential in the general form $V'(x)=\sum_{j=0}^{\deg V'}g_{j+1}x^j$. First,
	\begin{align}
		V'(x_k)W(x_k)-R(x_k) &= \sum_{n=0}^\infty\sum_{j=0}^n \frac{g_{j+1}\mu_n}{x_{k}^{n+1-j}} = \underbrace{\sum_{n=0}^\infty \sum_{j=0}^{n-1} \frac{g_{j+1}\mu_n}{x_k^{n+1-j}}}_{A}+\underbrace{\sum_{n=0}^{\deg V'} \frac{g_{n+1}\mu_n}{x_k}}_{B} \label{cotdecomp}
	\end{align}
	The term $A$ may be bounded in a similar manner,
	\begin{align*}
		|\text{(a)}| &\leq \sum_{j=0}^{\deg V'}\frac{1}{|x_k|^{1-j}}\sum_{n=j+1}^\infty \frac{g_{j+1}L^n}{|x_k|^{n}} = \sum_{j=0}^{\deg V'} \frac{g_{j+1}L^{j+1}}{|x_k|(|x_k|-L)}\leq G_1 \frac{1}{(|k|m-R)(|k|m-R-L)}
	\end{align*}
	where $G_1 = \sum_{j=0}^{\deg V'} g_{j+1}L^{j+1}$. For the term $B$, set $G_2 = \sum_{n=0}^{\deg V'}g_{n+1}\mu_n$, so that
	\[ B = G_2 \frac{1}{x_k}. \] 
	
	Our next step is to apply the M-test to $H(x)$ on the ball $B(0;R)$. Set $N=\lceil(R+L)/m\rceil$. We split $S$ into two terms,
	\[ H(x) = \sum_{|k|\leq N} f_k(x) + \sum_{|k|\geq N} f_k(x) \]
	The first term may only have cuts at $I_{\pm N}$ as previously discussed, and $I_{\pm N} = \Sigma \pm \lceil R+L\rceil i$, thus the cuts lie outside of $B(0;R)$, hence the first term $\sum_{|k|\leq N} f_k(x)$ is analytic on $B(0;R)$. For the second term, by our choice of $N$, we may expand each $f_k$ and use our previous estimates to find 
	\begin{align}
		|f_k(x)| &\leq |W(x_k)|^2 + \frac{\beta_2}{2}|W(x_k)||W(x_{k+1})| + \frac{4}{\beta}|V'(x_k)W(x_k)-R(x_k)| \nonumber\\
		&\leq \frac{1}{(|k|m-(R+L))^2}+\frac{\beta_2}{2}\frac{1}{(|k|m-(R+L))(|k+1|m-(R+L))}\nonumber\\&\qquad+\frac{4}{\beta}G_1\frac{1}{(|k|m-R)(|k|m-R-L)} + \frac{4}{\beta}G_2 \frac{1}{|x_k|}. \label{fbound}
	\end{align} 
	Upon summing over $k$, all the terms, except for the last one, clearly converge. However, note that
	\[ \sum_{k\in\zz} \frac{1}{x_k} = \frac{1}{mi}\sum_{k\in\zz} \frac{1}{x/mi-k} = \frac{\pi}{mi}\cot\frac{\pi x}{mi} \]
	by the well known partial fraction expansion of $\cot$. Hence $\sum_{|k|\geq N} \frac{1}{x_k}$ is analytic on $B(0;R)$. Thus, we may apply the M-test to $\sum_{k\in\zz} f_k$ for any $B(0;R)$, hence showing that $\sum_{k\in\zz} f_k$ uniformly converges to a continuous function on any compact subset of $\cc$, and therefore to a holomorphic function on $\cc$ by an application of Morera's theorem.
\end{proof}

\end{document}
